%/*% pdflatex -file-line-error -interaction nonstopmode %

\documentclass{article}
\usepackage{graphicx}
\usepackage{amsmath}
\usepackage{amsfonts}
\usepackage{bookmark}
\usepackage{booktabs}
\usepackage{exercises}
\graphicspath{ {./images/test} }
\usepackage{indentfirst}
\usepackage{amsthm}
\usepackage{yhmath}
\usepackage{placeins}
\newcommand*{\oldneg}{\mathord{\sim}}
\newcommand\showdiv[1]{\overline{\smash{)}#1}}
\newcommand\Mydiv[1]{%
\smash{\raise.3ex\hbox{$\big)$}}\mkern-8mu
        \overline{\enspace\strut#1}}
\setlength{\tabcolsep}{5pt}
\theoremstyle{definition}
\newtheorem{definition}{Definition}

\setlength{\parindent}{.5in}
\setlength{\parskip}{12pt}
\renewcommand{\arraystretch}{2}

\title{Quiz 1 (pg. 179)}
\begin{document}
\maketitle

\section*{1}
\[ 36^\frac{3}{2}=(\sqrt{36})^3=6^3 \]

\section*{2}
\[ 64^{-\frac{2}{3}}=(\sqrt[3]{64})^{-2}=4^{-2}=\frac{1}{4^3}=\frac{1}{16} \]

\section*{3}
\[ -(625^\frac{3}{4})=-(\sqrt[4]{625})^3=-(5^3)=-125 \]

\section*{4} 
\[ (-32)^\frac{2}{5}=(\sqrt[5]{-32})^2=(-2)^2=4 \]

\section*{5}
\[ x^4=20 \] \[ x=\pm\sqrt[4]{20}=\pm 2.11 \]

\section*{6}
\[ x^5=-10 \] \[ x=\sqrt[5]{-10}=-1.58 \]

\section*{7}
\[ x^6+5=26 \] \[ x=\pm\sqrt[6]{21}=\pm 1.66 \]

\section*{8}
\[ (x+3)^3=-16 \] \[ x=-3+\sqrt[3]{-16}=-5.52 \]

\section*{9}
\[ \sqrt[4]{32}\times\sqrt[4]{8}=2\sqrt[4]{2}\times\sqrt[4]{8}=2\sqrt[4]{16}=2\times 2=4 \]

\section*{10}
\[ (\sqrt{10}\times \sqrt[3]{10})^8=10^4\times \sqrt[3]{10}^8=10^\frac{20}{3}=10^6\sqrt[3]{100} \]

\section*{11}
\[ (x^6y^4)^\frac{1}{8}+2(x^\frac{1}{3}y^\frac{1}{4})^2=x^\frac{3}{4}y^\frac{1}{2}+2x^\frac{2}{3}y^\frac{1}{2} \]

\section*{12}
\[ \frac{3\sqrt{7^3}+4\sqrt{7^3}}{\sqrt{7^5}}=\frac{7\sqrt{7^3}}{\sqrt{7^5}}=\frac{49\sqrt{7}}{\sqrt{7^5}} \]
\[ =\frac{49\sqrt{7}}{49\sqrt{7}}=1 \]

\section*{13}
\[ \frac{2\sqrt{x}\times\sqrt{x^3}}{\sqrt{64x^{15}}}=\frac{2\sqrt{x}\times x\sqrt{x}}{8x^7\sqrt{x}} \]
\[ =\frac{\sqrt{x}}{4x^6} \]

\section*{14}
\[ y^2\sqrt[5]{64x^6}-6\sqrt[5]{2x^6y^{10}}=2y^2x\sqrt[5]{2x}-6y^2x\sqrt[5]{2x} \]
\[ =-4y^2x\sqrt[5]{2x} \]

\section*{15} 
First, we can show that the inscribed red triangle is a right triangle. Because the two triangles with side lengths shown are both right triangles 
and have two proportional sides, we know that these two triangles are similar. Then, look at the three angles forming the straight angle on the side of 
the square that is split into two parts of 4 units. Call the part of this straight angle that is part of the red triangle $\angle{A}$, the one that is 
part of the bigger labeled triangle $\angle{B}$, and the other angle (part of the smaller labeled triangle) $\angle{C}$. Thus 
\[ \angle{A}+\angle{B}+\angle{C}=180 \] Because the two labeled triangles are similar and right, we can deduce that $\angle{B}+\angle{C}=90$, as 
$\angle{B}$ and $\angle{C}$ are not equal (look at where the corresponding sides are; these two angles have different corresponding sides in the other 
triangle). Thus \[ \angle{A}+90=180 \] thus confirming that the red triangle is a right triangle. 
We know that one of the sides of the red triangle is $\sqrt{8^2+4^2}=4\sqrt{5}$, and another side of the right triangle is 
$\sqrt{2^2+4^2}=\sqrt{20}=2\sqrt{5}$. By the pythagorean theorem, the last side of the red triangle is 
$\sqrt{(2\sqrt{5})^2+(4\sqrt{5})^2}=10$, meaning the perimeter of the red triangle is 
\[ P=10+4\sqrt{5}+2\sqrt{5}=10+6\sqrt{5} \] units.

\end{document}
